\section{System and Agent Workflow}
\label{sec:system}

\subsection{Interfaces and Shared Data Model}
Figure~\ref{fig:arch} summarizes the platform. The Nuxt web console
supports collection browsing, document management, preview, parsing, and
search; a FastAPI backend schedules MinerU parsing jobs~\cite{mineru}.
The \texttt{ragctl} CLI and MCP server expose the same resources to
operators and agents. The current MCP inventory contains 42
\texttt{kb\_*} tools within 95 tools; packaged skills combine them into
reusable ingest, search, management, graph, experience, and persona
workflows. Skills add workflow policy, while the clients share the same
underlying collections.

Each collection is a hierarchy of bases, documents, parts, and sections.
Markdown files are paired with two inspectable registries:
\path{.tree-fs.json} stores the tree, while \path{.knowledge-base.yml}
stores document descriptions, tags, source addresses, and index state.
ChromaDB stores BGE-M3 vectors~\cite{bge-m3};
an optional Neo4j index stores document relations. The file-backed
registries remain inspectable, and derived indexes are checked against
their recorded base identifiers.

\subsection{Gated Ingestion}
The ingest workflow converts PDFs to structured Markdown, checks parse
quality, detects exact and near duplicates, and splits oversized files
at structural boundaries (30k--60k characters, depending on
configuration). Content analysis routes each document or part to an
appropriate base or sub-base. The workflow adds 2--5 normalized tags and
a description covering five query dimensions: domain, method, object,
answerable problem, and conclusion. Each description is checked against
the stored source. Indexing follows supported write paths automatically.
The final C1--C9 checklist checks completeness, description and tag
quality, base attribution, vector and graph indexing, parsed images,
registry-to-disk consistency, and retrieval through the description.
An indexing or consistency failure blocks completion.

Table~\ref{tab:pipeline} reports a saved audit of the 100-paper
demonstration collection. The audit sampled the first 1{,}000 characters
of each part for its content check; it is not a full semantic audit of
every character.

\begin{table}[t]
\caption{Ingest audit on the 100-paper demonstration collection. Values
are script-asserted from saved artifacts. The content check samples the
first 1{,}000 characters per part; storage checks cover 33 units.}
\label{tab:pipeline}
\footnotesize
\setlength{\tabcolsep}{3.5pt}
\begin{tabular}{@{}lc@{}}
\toprule
\textbf{Stage / check} & \textbf{Result} \\
\midrule
Parsed and content-routed papers & 100/100 across 5 bases \\
Stored parts after the 30k split ceiling & 322 \\
Per-paper C1--C9 checks & 100/100 pass \\
True duplicates & 0 (35 near-pairs flagged) \\
Descriptions upgraded and read back & 165/165 \\
Part-content / storage audits & 322/322; 33/33 \\
Retrieval regression probes & 15/15 \\
\bottomrule
\end{tabular}
\end{table}

The same registries support document and collection CRUD, tag
administration, optional relation-graph construction, and reusable
experience entries extracted from resolved cases. The MCP inventory
also exposes SOUL persona operations; these are available in the
platform but are outside the retrieval ablation below.

\subsection{Layered Retrieval and Evidence Judgment}
\label{sec:ladder}
The retrieval strategy, \sys{} (query-driven, content-verified retrieval),
uses the shared address model and separates candidate recall from
evidence acceptance. As Figure~\ref{fig:arch} shows, the workflow has
six steps:
\emph{L0, catalog}. A lightweight call reads base names, descriptions,
and document counts. A base is pruned only when its description is
present and clearly irrelevant.
\emph{L1, descriptions}. The agent reads document cards in the remaining
bases; descriptions narrow the search but do not establish evidence.
\emph{L2, recall}. The query is rewritten as a declarative statement
with key terms. A wide vector search (top-$k=30$, threshold 0.35, with
one retry at 0.30) proposes candidates; literal enumeration questions
also use an exact-match scan.
\emph{L3, judgment}. The local multilingual Laya engine scores each
candidate's 3{,}000-character segments for answerable evidence. A
segment is retained when its positive-evidence probability
(\texttt{noul}) is at least 0.5. Topical similarity alone is
insufficient. Missing or invalid scores fail closed, so the lane never
falls back to similarity-only acceptance.
\emph{L4, reading}. The agent reads surviving evidence by its recorded
base/document/part/section address and checks key values against the
source.
\emph{L5, answer}. The answer cites the inspected source; the trace
records the search path, retrieval backend, gate threshold, and
blind spots. The returned evidence pack is subject to a configured
context cap.

If the catalog and descriptions reveal no plausible source, or an exact
scan finds no requested item, the workflow returns a scoped not-found
report. It records the searched bases and retrieval stages rather than
continuing with unbounded retries.

\subsection{Retrieval Lanes}
\sys{} defines a shared contract; three lanes trade breadth and latency.
The \emph{fast lane} rewrites the query, runs wide vector recall and a
BM25-style cross-check, then sends deduplicated candidates to the same
evidence gate. The \emph{librarian lane} navigates base and document
descriptions and uses exact scans when the relevant shelf is unknown or
complete recall is requested; it does not use vector recall. The
\emph{fusion lane} runs both lanes concurrently and merges candidates
by normalized base/document identity before judgment. All lanes return
source references and evidence notes to the answering agent.
